\section{5 · The Pratyakṣa System: Architecture}\label{the-pratyakux1e63a-system-architecture}

We describe the system as a layered architecture: a Core ContextStore, three runtime modules (Avacchedaka query, Sublation, Bayesian aggregation), three reasoning modules (Buddhi/Manas, Sākṣī, Khyātivāda classifier), and three resource-management modules (EventBoundaryCompactor, AdaptiveForgetting, TokenBudgetWatchdog). Each module is implemented as one or more MCP tools (Section 6.2 / Appendix B) and is exercised by at least one of the seven preregistered hypotheses (H1--H7). Figure \ref{fig:architecture} shows the deployed surface end-to-end: 15 MCP tools (grouped into four families) consumed by 3 sub-agents, exposed through 4 slash-commands and 3 lifecycle hooks, and mounted by all four host surfaces (Claude Code CLI, Claude Code VS Code, Cursor, Claude Desktop) without modification of the underlying LLM.

% =====================================================================
% Fig 1 -- Harness architecture overview
% =====================================================================
\begin{figure}[!tbp]
\centering
\resizebox{\linewidth}{!}{%
\begin{tikzpicture}[
  font=\small,
  every node/.style={rounded corners=2pt},
  fam/.style ={draw=black!60,fill=blue!6,inner sep=4pt,minimum height=10mm,
               text width=27mm,align=center},
  agent/.style={draw=black!60,fill=orange!8,inner sep=3pt,minimum height=8mm,
               text width=24mm,align=center},
  cmd/.style ={draw=black!60,fill=green!7,inner sep=3pt,minimum height=8mm,
               text width=22mm,align=center},
  hook/.style={draw=black!60,fill=violet!8,inner sep=3pt,minimum height=8mm,
               text width=22mm,align=center},
  host/.style={draw=black!60,fill=gray!10,inner sep=4pt,minimum height=12mm,
               text width=30mm,align=center},
  arr/.style ={-{Latex[length=2mm]},gray!70,thick},
  >=Latex
]
% --- MCP tool families (left column) ---
\node[fam] (insert)   {\textbf{Context store (5)}\\
  context\_insert,\\ context\_retrieve,\\ context\_window,\\ list\_items, get\_item};
\node[fam,below=4mm of insert] (sublate) {\textbf{Sublation (3)}\\
  detect\_conflict,\\ sublate\_with\_evidence,\\ list\_sublations};
\node[fam,below=4mm of sublate] (witness){\textbf{Witness/audit (3)}\\
  set\_sakshi,\\ get\_sakshi\_log,\\ status\_snapshot};
\node[fam,below=4mm of witness] (judge)  {\textbf{Judgment (4)}\\
  classify\_khyativada,\\ token\_budget,\\ adaptive\_forget,\\ event\_boundary};

% --- Agents (middle column) ---
\node[agent,right=18mm of insert] (manas)  {\textbf{ManasAgent}\\ select \& attend};
\node[agent,below=3mm of manas]   (buddhi) {\textbf{BuddhiAgent}\\ judge \& tag};
\node[agent,below=3mm of buddhi]  (sakshi) {\textbf{SakshiKeeper}\\ write-once log};

% --- Commands & hooks ---
\node[cmd,right=12mm of manas]  (cmd1) {/context-status};
\node[cmd,below=2mm of cmd1]    (cmd2) {/context-doctor};
\node[cmd,below=2mm of cmd2]    (cmd3) {/sublate};
\node[cmd,below=2mm of cmd3]    (cmd4) {/witness};

\node[hook,below=8mm of cmd4]  (h1) {pre-tool hook};
\node[hook,below=2mm of h1]    (h2) {post-tool hook};
\node[hook,below=2mm of h2]    (h3) {session hook};

% --- Hosts (right column) ---
\node[host,right=14mm of cmd2] (cli)   {Claude Code\\ \textit{CLI}};
\node[host,below=3mm of cli]   (vsc)   {Claude Code\\ \textit{VS Code}};
\node[host,below=3mm of vsc]   (cur)   {\textit{Cursor}};
\node[host,below=3mm of cur]   (desk)  {Claude\\ \textit{Desktop}};

% --- Arrows: tools -> agents ---
\draw[arr] (insert.east)  -- (manas.west);
\draw[arr] (sublate.east) -- (buddhi.west);
\draw[arr] (witness.east) -- (sakshi.west);
\draw[arr] (judge.east)   -- (buddhi.west);

% --- Arrows: agents <-> commands ---
\draw[arr] (manas.east)  -- (cmd1.west);
\draw[arr] (buddhi.east) -- (cmd2.west);
\draw[arr] (sakshi.east) -- (cmd4.west);

% --- Arrows: hooks observe agents ---
\draw[arr,dashed] (h1.west) -- ++(-3mm,0) |- (manas.east);
\draw[arr,dashed] (h2.west) -- ++(-3mm,0) |- (buddhi.east);
\draw[arr,dashed] (h3.west) -- ++(-3mm,0) |- (sakshi.east);

% --- Arrows: hosts mount everything ---
\draw[arr] (cmd1.east) -- (cli.west);
\draw[arr] (cmd2.east) -- (vsc.west);
\draw[arr] (cmd3.east) -- (cur.west);
\draw[arr] (cmd4.east) -- (desk.west);

% --- Frames ---
\begin{pgfonlayer}{background}
  \node[draw=black!30,dashed,inner sep=4mm,fit=(insert)(judge),
        label={[font=\footnotesize\itshape]above:15 MCP tools (server)}]{};
  \node[draw=black!30,dashed,inner sep=4mm,fit=(manas)(sakshi),
        label={[font=\footnotesize\itshape]above:3 sub-agents}]{};
  \node[draw=black!30,dashed,inner sep=4mm,fit=(cmd1)(h3),
        label={[font=\footnotesize\itshape]above:4 slash cmds + 3 hooks}]{};
  \node[draw=black!30,dashed,inner sep=4mm,fit=(cli)(desk),
        label={[font=\footnotesize\itshape]above:Host surfaces}]{};
\end{pgfonlayer}
\end{tikzpicture}}
\caption{\textbf{Pratyak\d{s}a system architecture.}
\textit{Fifteen MCP tools} (grouped into four families: context store, sublation, witness/audit, and judgment) are consumed by \textit{three sub-agents} (\textbf{Manas} attends and selects; \textbf{Buddhi} judges and emits; \textbf{Sak\d{s}\=\i Keeper} writes the immutable per-turn audit log).
\textit{Four slash-commands} and \textit{three lifecycle hooks} expose this machinery to four host surfaces (Claude Code CLI, Claude Code VS Code, Cursor, and Claude Desktop) without modification of the underlying LLM.
Solid arrows are runtime calls; dashed arrows are observation/audit.}
\label{fig:architecture}
\end{figure}


\subsection{5.1 The ContextStore (core)}\label{the-contextstore-core}

The \passthrough{\lstinline!ContextStore!} is a typed, append-mostly key-value store that lives outside the LLM and is exposed to it via MCP. Its row schema is exactly the \emph{avacchedaka} tuple of Section 4.2:

\begin{lstlisting}[language=Python]
@dataclass(frozen=True)
class ContextItem:
    id: str
    qualificand: str         # the "thing" being asserted about
    qualifier: str           # what is asserted about it
    condition: str           # the limitor under which it holds
    precision: float         # source-supplied calibrated reliability ∈ [0, 1]
    source: str              # "user" | "rag" | "tool:foo" | "agent:buddhi" | ...
    inserted_at: float       # monotonic seconds
    stale: bool              # source-flagged staleness
    status: Literal["live", "bādhita", "evicted"]
    superseded_by_id: str | None
    witness_id: str          # set by SakshiKeeperAgent on first witness
\end{lstlisting}

Insertion is via a single MCP tool, \passthrough{\lstinline!context\_insert(...)!}. Retrieval is via \passthrough{\lstinline!context\_retrieve(...)!} (with optional \passthrough{\lstinline!qualificand!}, \passthrough{\lstinline!qualifier\_substring!}, and \passthrough{\lstinline!condition!} filters) and \passthrough{\lstinline!context\_get(...)!} for direct id fetch; none of these return \passthrough{\lstinline!bādhita!} or \passthrough{\lstinline!evicted!} rows by default. A debug flag exposes them for audit.

Two design commitments distinguish this from a vector store:

\begin{enumerate}
\def\labelenumi{\arabic{enumi}.}
\item
  \textbf{Conditions are first-class.} Two rows that share \passthrough{\lstinline!qualificand!} and disagree on \passthrough{\lstinline!qualifier!} are \emph{not} yet in conflict if their \passthrough{\lstinline!condition!} fields differ. This is the avacchedaka commitment.
\item
  \textbf{Provenance never disappears.} Sublation rewrites \passthrough{\lstinline!status!}, never deletes. Eviction (under hard budget pressure) writes to a sidecar \passthrough{\lstinline!evicted\_log.jsonl!} so the witness can replay the eviction reason.
\end{enumerate}

\subsection{5.2 The AvacchedakaQuery module}\label{the-avacchedakaquery-module}

\passthrough{\lstinline!context\_retrieve!} with \passthrough{\lstinline!qualificand!} and \passthrough{\lstinline!condition!} set walks the \passthrough{\lstinline!ContextStore!} and returns the \emph{set} of items matching \passthrough{\lstinline!qualificand!} whose \passthrough{\lstinline!condition!} is \emph{consistent} with the supplied condition. Consistency is decided by a small rule engine:

\begin{itemize}
\tightlist
\item
  \textbf{String equality} wins when conditions are identical.
\item
  \textbf{Temporal subsumption} wins when one condition is a strictly tighter time-window of the other (e.g.~\passthrough{\lstinline!"Django ≥ 5.0"!} subsumes \passthrough{\lstinline!"Django 5.0.4"!}).
\item
  \textbf{Lexical-prefix subsumption} for versioned platforms (\passthrough{\lstinline!"POSIX"!} subsumes \passthrough{\lstinline!"POSIX, Python ≥ 3.11"!}).
\item
  \textbf{Otherwise}, items are returned in \emph{both} groups and the conflict is left to be resolved downstream by Buddhi, possibly via \passthrough{\lstinline!sublate\_with\_evidence!}.
\end{itemize}

This rule engine is intentionally simple. We are not building a description-logic reasoner; we are providing the agent with a typed surface on which to act.

\subsection{5.3 The Sublation + Bayesian aggregation module}\label{the-sublation-bayesian-aggregation-module}

Conflict resolution proceeds in two steps. First, \passthrough{\lstinline!sublate\_with\_evidence(target\_id, by\_id, reason)!} performs the Vedānta operation of Section 4.3. Second, when the agent must commit to a single posterior probability for a contested claim --- for example, ``is the user's setting \passthrough{\lstinline!enable\_x = True!}?'' --- the harness offers a \textbf{Bayesian Beta-Bernoulli aggregator} in place of the original \emph{PrecisionWeightedRAG} of v0:

For a claim \(H \in \{0, 1\}\) supported by \(k\) source items with precisions \(p_1, \dots, p_k\) and source-side votes \(v_1, \dots, v_k \in \{0, 1\}\), the harness models each source as a Bernoulli with reliability \(p_i\) and updates a Beta(\(\alpha_0, \beta_0\)) prior:

\[
\alpha_n = \alpha_0 + \sum_{i: v_i=1} p_i,\qquad \beta_n = \beta_0 + \sum_{i: v_i=0} p_i.
\]

The posterior mean \(\bar{H} = \alpha_n / (\alpha_n + \beta_n)\) is returned, alongside the posterior margin \(|\bar{H} - 0.5|\). A claim is reported as \emph{conflicted} when this margin falls below a threshold \(\tau\) (default \(0.10\)). Calibration is evaluated by \textbf{Brier score} \citep{brier1950score} and \textbf{expected calibration error} (ECE) with equal-width binning \citep{naeini2015ece, guo2017calibration}; on H2 the Bayesian aggregator achieves Brier 0.094 vs.~PrecisionWeightedRAG 0.176 and ECE 0.041 vs.~0.118 on a held-out 1,800-example validation slice (Section 8.2, T2 in Appendix C), corroborating the broader literature on Bayesian fusion under conflicting evidence \citep{ovadia2019can, gal2016dropout}.

\subsection{5.4 Buddhi and Manas as plug-in sub-agents}\label{buddhi-and-manas-as-plug-in-sub-agents}

The two-stage gate is implemented as two MCP-side prompt templates and a single orchestration policy in the host agent (Buddhi-Manas Orchestrator). The data-flow is:

\begin{enumerate}
\def\labelenumi{\arabic{enumi}.}
\tightlist
\item
  The host agent enters the harness's \passthrough{\lstinline!manas\_step(query)!} which returns the Manas JSON.
\item
  The host agent enters \passthrough{\lstinline!buddhi\_step(query, manas\_output)!} which returns the final answer + khyāti class.
\item
  The host agent enters \passthrough{\lstinline!sakshi\_record(...)!} which appends the immutable witness log entry.
\end{enumerate}

Both Manas and Buddhi are \emph{prompted-only} sub-agents. We deliberately do \emph{not} fine-tune. The two-stage gate is a discipline, not a model: a different LLM (Claude 3.5 Sonnet, Claude 4.x, GPT-4o, Qwen-3) can be plugged in as the \emph{runner} of either prompt without breaking the harness's invariants. This is the ``host-platform-agnostic'' design commitment, and it makes the plugin hot-swappable across Cursor / Claude Code / Claude Desktop \citep{anthropic2025mcp, cursor2025plugins}.

\subsection{5.5 The Khyātivāda classifier}\label{the-khyux101tivux101da-classifier}

The shipped \passthrough{\lstinline!classify\_khyativada!} tool applies a \textbf{heuristic labeler with structured JSON shape} plus the rule-based guardrails below. The experiments harness additionally implements a few-shot Anthropic JSON variant (\passthrough{\lstinline!src/evaluation/khyativada\_fewshot.py!}) for offline evaluation; that path is \textbf{not} wired into the plugin MCP surface. For exposition we still describe the JSON contract shared by both stacks. Inputs:

\begin{itemize}
\tightlist
\item
  The user query.
\item
  The Buddhi answer.
\item
  (Optional) the Manas-attended items.
\end{itemize}

Output:

\begin{lstlisting}
{
  "khyati_class": "anyathākhyāti" | "ātmakhyāti" | ... | "none",
  "rationale": "<≤2 sentence explanation citing item-ids and qualifiers>",
  "confidence": <float in [0,1]>
}
\end{lstlisting}

Sample exemplars are curated for the few-shot \emph{experiment} path to maximize within-class lexical diversity (e.g., asatkhyāti includes both ``fabricated CVE'' and ``fabricated stack-trace line'', not just one). In both stacks, heuristic guardrails override the model-proposed label when (a) the answer cites a \passthrough{\lstinline!qualificand!} not present in any \passthrough{\lstinline!ContextStore!} item (force \passthrough{\lstinline!asatkhyāti!}), or (b) the answer's \passthrough{\lstinline!qualifier!} is the negation of an item the answer cites by id (force \passthrough{\lstinline!viparītakhyāti!}). These guardrails are \emph{additive} --- they only flip \passthrough{\lstinline!none!} to a class, never flip one class to another.

\subsection{5.6 The EventBoundaryCompactor}\label{the-eventboundarycompactor}

For long-running sessions we want to compact whole \emph{episodes} of context, not individual items. Following event-segmentation theory \citep{zacks2007event, baldassano2017nested} and predictive-coding accounts of surprise \citep{rao1999predictive, friston2010fep, feldman2010precision}, we segment the rolling token stream by per-token \textbf{negative-log-probability surprise} computed by a small open-weights model --- Qwen3-1.7B served via vLLM \citep{kwon2023vllm, vllm2023, qwen2024technical}, with an optional fallback to Qwen3-0.6B and a final heuristic-Zipf fallback \citep{piantadosi2014zipf, shannon1948mathematical, mackay2003information} when the GPU path is unavailable.

A boundary is declared when surprise exceeds a rolling-mean+2σ threshold sustained across \(w=12\) tokens. At boundary, the EventBoundaryCompactor:

\begin{enumerate}
\def\labelenumi{\arabic{enumi}.}
\tightlist
\item
  Summarises the just-closed episode into a single new \passthrough{\lstinline!ContextItem!} with \passthrough{\lstinline!qualifier="episode\_summary"!} and \passthrough{\lstinline!precision = mean(precisions of contained items)!}.
\item
  Marks all contained items as \passthrough{\lstinline!status="evicted"!} (provenance retained).
\item
  Notifies the SakshiKeeperAgent of the boundary.
\end{enumerate}

H4 (Section 8.4) tests this module against a no-compaction baseline.

\subsection{5.7 AdaptiveForgetting}\label{adaptiveforgetting}

Rules are as in Section 4.7. Implemented as a periodic background MCP tool \passthrough{\lstinline!compact(strategy: Literal["adaptive", "lru", "none"])!} that the agent or a lifecycle hook can invoke. The default policy is \passthrough{\lstinline!adaptive!}; the H7 study (Section 8.7) compares all three.

\subsection{5.8 TokenBudgetWatchdog}\label{tokenbudgetwatchdog}

A small but indispensable module. The watchdog tracks input/output tokens by category (system, retrieval, tool, completion) against a configured budget (default 60\% of model context window) and exposes one MCP tool, \passthrough{\lstinline!budget\_status()!}, returning:

\begin{lstlisting}
{
  "context_tokens_used": 12345,
  "context_budget": 16384,
  "fraction_used": 0.753,
  "by_category": {"system": 1024, "retrieval": 7000, "tool": 3000, "completion": 1321},
  "compaction_recommended": false
}
\end{lstlisting}

It is also wired into a \passthrough{\lstinline!pretooluse-budget.sh!} hook (Section 6.4) that \textbf{logs} budget pressure by default; optional strict refusal is available when \passthrough{\lstinline!PRATYAKSHA\_BUDGET\_STRICT=1!}.

\subsection{5.9 The full request lifecycle}\label{the-full-request-lifecycle}

A representative single-turn request now flows as:

\begin{enumerate}
\def\labelenumi{\arabic{enumi}.}
\tightlist
\item
  \textbf{User query} lands in the host agent (Claude Code, Cursor, or Claude Desktop).
\item
  \textbf{Manas} is called: it inspects the \passthrough{\lstinline!ContextStore!} snapshot via \passthrough{\lstinline!context\_retrieve!} / \passthrough{\lstinline!context\_window!} and returns its attended subset.
\item
  \textbf{Buddhi} is called with that subset. It may issue \textbf{\passthrough{\lstinline!sublate\_with\_evidence!}} for any conflicts it spots before answering.
\item
  \textbf{\passthrough{\lstinline!classify\_khyativada!}} is run on Buddhi's draft answer; if \passthrough{\lstinline'khyati\_class != "none"'} and \passthrough{\lstinline!confidence > τ\_self\_check!} (default 0.7), Buddhi is asked to revise once.
\item
  \textbf{Sākṣī} writes the witness log entry.
\item
  The answer is returned to the user.
\item
  \textbf{AdaptiveForgetting} and \textbf{EventBoundaryCompactor} are invoked by lifecycle hooks if budget pressure is detected.
\end{enumerate}

This lifecycle is the running invariant tested by every experiment in Sections 8--10.
